% Documento:    Contenido del informe - Laboratorio Kafka (rehecho) + Laboratorio 07 Flink
% Integrante:   Jorge Tito Ccahuaya
% Curso:        Big Data (1705157)
%
% Este documento está pensado para funcionar también como guía reproducible:
% cada evidencia (captura) va precedida del comando exacto y copiable que la
% generó.

% ------------------------------------------------------------------------------
\section{Introducción}
% ------------------------------------------------------------------------------

	Este informe documenta dos laboratorios del curso Big Data desarrollados sobre una misma infraestructura distribuida en AWS: la reconstrucción completa, desde cero, del laboratorio de Apache Kafka (clúster de 3 brokers sobre instancias EC2, demostrando productor, consumidor, tópicos, particiones, offsets y grupos de consumidores), y el nuevo Laboratorio 07 de Apache Flink, que consume los eventos generados por Kafka para aplicar filtrado, transformación, conteo y agrupamiento en tiempo real.

	A diferencia de un informe puramente descriptivo, este documento se escribió para que también sirva de guía: cada paso relevante muestra primero el comando exactamente como se ejecutó, y luego la evidencia (captura de pantalla o salida de consola) de su resultado, de modo que el procedimiento completo pueda reproducirse sin ambigüedad.

	Los fragmentos de código mostrados a lo largo del informe (código de infraestructura en Pulumi, scripts de Kafka, productor, consumidor y el proyecto de Flink) son extractos ilustrativos; el código fuente completo y actualizado está disponible en el repositorio del proyecto: \href{https://github.com/Noodle96/Big_Data/tree/main/kafka-flink-streaming-lab}{github.com/Noodle96/Big\_Data/tree/main/kafka-flink-streaming-lab}.


% ------------------------------------------------------------------------------
\section{Infraestructura del clúster Kafka}
% ------------------------------------------------------------------------------

	\subsection{Arquitectura desplegada}

		La infraestructura se aprovisiona con Pulumi (Python) sobre AWS e-2. Se creó una VPC dedicada (\texttt{10.20.0.0/16}) con una subred pública (\texttt{10.20.1.0/24}), donde conviven cuatro instancias EC2 Ubuntu 24.04 (\texttt{t3.small}): tres brokers de Kafka (\texttt{kafka-broker-1/2/3}) y un cliente (\texttt{kafka-client}) desde donde se ejecutan el productor y el consumidor. Todas comparten un mismo Security Group que permite tráfico interno total entre ellas y acceso SSH desde Internet. Cada instancia tiene una IP privada fija asignada de antemano (brokers en \texttt{10.20.1.11-13}, cliente en \texttt{10.20.1.20}), lo que permite generar de forma determinística la configuración de cada nodo antes incluso de crear las instancias.

	\subsection{Aprovisionamiento automático (\textit{user data})}

		Cada instancia EC2 recibe, en el momento de su creación, un script (\textit{user data}) que AWS ejecuta automáticamente mediante \texttt{cloud-init} apenas la máquina arranca por primera vez. Este mecanismo es el que garantiza la reproducibilidad del laboratorio: ningún broker se configuró manualmente por SSH; toda la instalación y configuración de Kafka ocurre de forma automática al momento del despliegue.

		\begin{sourcecode}[\label{codigo-pulumi-up}]{bash}{Despliegue de toda la infraestructura con Pulumi.}
cd infrastructure
pulumi up
		\end{sourcecode}

		La Figura \ref{img:pulumi-up} muestra la salida real de este comando, confirmando la creación de los 15 recursos declarados en el código (VPC, subred, Security Group, sus reglas, y las 4 instancias EC2).

		\setindexcaption{Resultado de \texttt{pulumi up}: los 15 recursos de la infraestructura creados.}
		\insertimage[\label{img:pulumi-up}]{kafka/infraestruture.png}{width=13cm}{Salida de \texttt{pulumi up}: VPC, subred, Security Group, sus reglas, y las 4 instancias EC2 (3 brokers + cliente) creadas correctamente.}

	\subsection{Configuración de cada broker (modo KRaft)}

		Para que un broker quede correctamente configurado dentro del clúster, el \textit{user data} realiza, en orden, los siguientes pasos:

		\begin{enumerate}
			\item \textbf{Instalación de Java}, dependencia de la JVM sobre la que corre Kafka.

			\begin{sourcecode}[\label{codigo-install-java}]{bash}{Instalación de Java 17.}
sudo apt-get update -y
sudo apt-get install -y openjdk-17-jdk-headless
			\end{sourcecode}

			\item \textbf{Descarga e instalación del binario de Kafka} (misma versión, 3.9.0, en los tres brokers).

			\begin{sourcecode}[\label{codigo-install-kafka}]{bash}{Descarga e instalación de Kafka.}
curl -fsSL https://archive.apache.org/dist/kafka/3.9.0/kafka_2.13-3.9.0.tgz -o kafka.tgz
tar -xzf kafka.tgz
ln -sfn /opt/kafka/kafka_2.13-3.9.0 /opt/kafka/current
			\end{sourcecode}

			\item \textbf{Usuario de sistema dedicado} y directorio de datos, para no correr Kafka como \texttt{root}.

			\begin{sourcecode}[\label{codigo-kafka-user}]{bash}{Usuario de sistema y directorio de datos.}
sudo useradd --system --home-dir /opt/kafka --shell /usr/sbin/nologin kafka
sudo mkdir -p /var/lib/kafka/data
sudo chown -R kafka:kafka /opt/kafka /var/lib/kafka
			\end{sourcecode}

			\item \textbf{Generación de \texttt{server.properties}} en modo KRaft combinado (\texttt{broker,controller}), con el \texttt{node.id} propio del broker (1, 2 o 3) y la lista completa de votantes del quorum (\texttt{controller.quorum.voters}), idéntica en los tres nodos.

			\begin{sourcecode}[\label{codigo-server-properties}]{plaintext}{server.properties generado para kafka-broker-1 (node.id=1).}
process.roles=broker,controller
node.id=1
controller.quorum.voters=1@10.20.1.11:9093,2@10.20.1.12:9093,3@10.20.1.13:9093

listeners=PLAINTEXT://10.20.1.11:9092,CONTROLLER://10.20.1.11:9093
advertised.listeners=PLAINTEXT://10.20.1.11:9092
controller.listener.names=CONTROLLER
inter.broker.listener.name=PLAINTEXT
listener.security.protocol.map=CONTROLLER:PLAINTEXT,PLAINTEXT:PLAINTEXT

log.dirs=/var/lib/kafka/data
num.partitions=3
default.replication.factor=3
			\end{sourcecode}

			\item \textbf{Formateo del almacenamiento} con \texttt{kafka-storage.sh}, usando el mismo \texttt{CLUSTER\_ID} en los tres brokers (si el ID no coincide, los nodos no pertenecen al mismo clúster).

			\begin{sourcecode}[\label{codigo-storage-format}]{bash}{Formateo del storage KRaft.}
sudo -u kafka /opt/kafka/current/bin/kafka-storage.sh format \
    -t X2k7a_8UT5i145WnsLy8qQ \
    -c /opt/kafka/current/config/kraft-server.properties \
    --ignore-formatted
			\end{sourcecode}

			\item \textbf{Registro como servicio \texttt{systemd}}, para que el proceso quede corriendo de forma persistente y se reinicie solo ante una falla.

			\begin{sourcecode}[\label{codigo-systemd}]{bash}{Habilitar y arrancar el servicio de Kafka.}
sudo systemctl daemon-reload
sudo systemctl enable kafka
sudo systemctl start kafka
			\end{sourcecode}
		\end{enumerate}

	\subsection{Verificación: broker activo}

		Una vez completado el arranque, se verificó en los \textbf{tres} brokers que el servicio quedara activo de forma estable (no reiniciando en bucle) y se revisó el registro de arranque de Kafka. Para eso, primero hay que conectarse por SSH a cada broker usando su IP pública (la entrega \texttt{pulumi up} como \texttt{sshBroker1}, \texttt{sshBroker2}, \texttt{sshBroker3} en sus outputs; cambia si se vuelve a desplegar):

		\begin{sourcecode}[\label{codigo-ssh-broker1}]{bash}{Conexión SSH a kafka-broker-1 (desde infrastructure/).}
ssh -i ../keys/kafka-lab-key.pem ubuntu@54.147.218.194
		\end{sourcecode}

		Ya dentro del broker, se corren los mismos dos comandos:

		\begin{sourcecode}[\label{codigo-check-status}]{bash}{Verificación del servicio (mismo comando en los 3 brokers).}
sudo systemctl status kafka
sudo journalctl -u kafka -n 30 --no-pager
		\end{sourcecode}

		La Figura \ref{img:broker1-status} confirma que \texttt{kafka-broker-1} quedó activo de forma estable, y la Figura \ref{img:broker1-journal} muestra su registro de arranque, sin excepciones.

		\setindexcaption{Servicio de Kafka activo en kafka-broker-1.}
		\insertimage[\label{img:broker1-status}]{kafka/broker1-systemctl-status.png}{width=13cm}{kafka-broker-1: activo de forma estable, sin reinicios.}

		\setindexcaption{Registro de arranque de Kafka en kafka-broker-1.}
		\insertimage[\label{img:broker1-journal}]{kafka/broker1-journalctl.png}{width=13cm}{kafka-broker-1: arranque limpio, sin excepciones.}

		Se repitió la misma verificación, conectándose por SSH a cada uno, en \texttt{kafka-broker-2} (Figuras \ref{img:broker2-status} y \ref{img:broker2-journal}) y en \texttt{kafka-broker-3} (Figuras \ref{img:broker3-status} y \ref{img:broker3-journal}), con idéntico resultado:

		\begin{sourcecode}[\label{codigo-ssh-broker23}]{bash}{Conexión SSH a kafka-broker-2 y kafka-broker-3.}
ssh -i ../keys/kafka-lab-key.pem ubuntu@52.90.232.221   # kafka-broker-2
ssh -i ../keys/kafka-lab-key.pem ubuntu@52.55.101.14    # kafka-broker-3
		\end{sourcecode}

		\setindexcaption{Servicio de Kafka activo en kafka-broker-2.}
		\insertimage[\label{img:broker2-status}]{kafka/broker2-systemctl-status.png}{width=13cm}{kafka-broker-2: activo de forma estable, sin reinicios.}

		\setindexcaption{Registro de arranque de Kafka en kafka-broker-2.}
		\insertimage[\label{img:broker2-journal}]{kafka/broker2-journalctl.png}{width=13cm}{kafka-broker-2: arranque limpio, sin excepciones.}

		\setindexcaption{Servicio de Kafka activo en kafka-broker-3.}
		\insertimage[\label{img:broker3-status}]{kafka/broker3-systemctl-status.png}{width=13cm}{kafka-broker-3: activo de forma estable, sin reinicios.}

		\setindexcaption{Registro de arranque de Kafka en kafka-broker-3.}
		\insertimage[\label{img:broker3-journal}]{kafka/broker3-journalctl.png}{width=13cm}{kafka-broker-3: arranque limpio, sin excepciones.}

	\subsection{Incidente durante el despliegue: KAFKA-18281}

		En el primer intento de despliegue, el formateo del storage (paso 5 de la sección anterior) falló con la siguiente excepción:

		\begin{sourcecode}[\label{codigo-bug-error}]{plaintext}{Error obtenido en el primer intento de formateo.}
Exception in thread "main" java.lang.IllegalArgumentException: requirement failed:
advertised.listeners cannot use the nonroutable meta-address 0.0.0.0.
Use a routable IP address.
		\end{sourcecode}

		La Figura \ref{img:bug-properties} muestra el \texttt{server.properties} generado en ese primer intento, donde se aprecia el origen del problema.

		\setindexcaption{server.properties del primer intento, con los listeners atados a 0.0.0.0.}
		\insertimage[\label{img:bug-properties}]{kafka/broker1-server-properties-bug.png}{width=13cm}{Configuración inicial: tanto \texttt{listeners} como el listener CONTROLLER quedaban atados a la dirección no enrutable 0.0.0.0.}

		La causa es un defecto conocido de Kafka 3.9.0, reportado como \textbf{KAFKA-18281}: Kafka valida incorrectamente el listener \texttt{CONTROLLER} exigiéndole una dirección enrutable, aun cuando dicho listener nunca se declaró en \texttt{advertised.listeners} (que es lo correcto). La corrección aplicada, sin depender de un parche de versión, fue atar tanto \texttt{listeners} como \texttt{advertised.listeners} a la IP privada específica de cada broker en vez de \texttt{0.0.0.0}. Como el script de aprovisionamiento se ejecuta con \texttt{set -e}, el fallo interrumpía el script antes de llegar a crear el servicio \texttt{systemd}; por eso, además de corregir la configuración, se destruyó y recreó toda la infraestructura (\texttt{pulumi destroy} seguido de \texttt{pulumi up}) para comprobar que el nuevo código funciona de punta a punta sin intervención manual.

	\subsection{Verificación del quorum KRaft}

		La activación del servicio en cada broker por separado no garantiza, por sí sola, que los tres nodos hayan formado un mismo quorum. Para confirmarlo se consultó el estado del quorum desde uno de los brokers (basta con estar conectado por SSH a cualquiera de los tres, como en el Código \ref{codigo-ssh-broker1}):

		\begin{sourcecode}[\label{codigo-quorum-status}]{bash}{Consulta del estado del quorum KRaft.}
kafka-metadata-quorum.sh \
    --bootstrap-server 10.20.1.11:9092,10.20.1.12:9092,10.20.1.13:9092 \
    describe --status
		\end{sourcecode}

		\setindexcaption{Estado del quorum KRaft del clúster de 3 brokers.}
		\insertimage[\label{img:quorum-status}]{kafka/quorum-status.png}{width=13cm}{Salida de \texttt{kafka-metadata-quorum.sh describe --status}: los tres node.id aparecen como votantes activos del quorum.}

		La salida obtenida en la Figura \ref{img:quorum-status} confirma que el clúster quedó correctamente configurado:

		\begin{itemize}
			\item \texttt{ClusterId} coincide exactamente con el \texttt{CLUSTER\_ID} fijado en el código de infraestructura (\texttt{X2k7a\_8UT5i145WnsLy8qQ}), es decir, los tres brokers efectivamente pertenecen al mismo clúster.
			\item \texttt{LeaderId: 3} indica que, de los tres controladores, \texttt{kafka-broker-3} ganó la elección y actúa como líder del quorum.
			\item \texttt{LeaderEpoch: 6} es el número de época de la elección actual; un valor mayor a 0 confirma que ya se completó al menos una elección de líder.
			\item \texttt{CurrentVoters} lista los tres \texttt{node.id} (1, 2 y 3) junto a su endpoint \texttt{CONTROLLER}, confirmando que los tres brokers participan activamente del quorum (ninguno quedó aislado).
			\item \texttt{CurrentObservers} vacío es esperado en este laboratorio: no hay controladores adicionales en modo solo observador.
		\end{itemize}

		Con este resultado se da por completada y verificada la configuración del clúster distribuido de Kafka.


% ------------------------------------------------------------------------------
\section{Uso del clúster: tópicos, productor y consumidor}
% ------------------------------------------------------------------------------

	Con el clúster ya configurado, esta sección demuestra los conceptos pedidos por el docente: \textit{topic}, \textit{partition}, \textit{producer}, \textit{consumer}, \textit{offset} y \textit{consumer group}. Todo se ejecuta desde \texttt{kafka-client}, la instancia dedicada a correr el productor y el consumidor (no los brokers).

	\subsection{Conexión al cliente y copia de los archivos}

		Primero hay que conectarse por SSH a \texttt{kafka-client} (IP pública entregada por \texttt{pulumi up} como \texttt{sshClient}):

		\begin{sourcecode}[\label{codigo-ssh-client}]{bash}{Conexión SSH al cliente.}
ssh -i keys/kafka-lab-key.pem ubuntu@54.90.122.174
		\end{sourcecode}

		El cliente no tiene el repositorio completo, así que los archivos necesarios se copian por \texttt{scp} antes de usarlos (este comando se ejecuta desde la máquina local, no dentro del cliente):

		\begin{sourcecode}[\label{codigo-scp-topics}]{bash}{Copia de topics.yaml y create\_topics.py al cliente.}
scp -i keys/kafka-lab-key.pem topics.yaml ubuntu@54.90.122.174:~/
scp -i keys/kafka-lab-key.pem scripts/kafka/create_topics.py ubuntu@54.90.122.174:~/
		\end{sourcecode}

	\subsection{Creación del topic (Topic y Partition)}

		El topic \texttt{ecommerce-events} se crea (y se describe automáticamente) con el script \texttt{create\_topics.py}, que lee la definición declarativa de \texttt{topics.yaml} en vez de recibir los parámetros sueltos por consola:

		\begin{sourcecode}[\label{codigo-create-topics}]{bash}{Creación del topic desde el cliente.}
python3 create_topics.py --bootstrap-server 10.20.1.11:9092,10.20.1.12:9092,10.20.1.13:9092
		\end{sourcecode}

		\setindexcaption{Creación y descripción del topic ecommerce-events.}
		\insertimage[\label{img:topic-create}]{kafka/topic-create-describe.png}{width=13cm}{Salida de \texttt{kafka-topics.sh --create} y \texttt{--describe} para el topic ecommerce-events: 3 particiones, factor de replicación 3.}

		La Figura \ref{img:topic-create} contiene toda la evidencia de los conceptos \textbf{Topic} y \textbf{Partition}. Interpretando cada campo de la salida:

		\begin{itemize}
			\item \texttt{TopicId} (\texttt{wz45FgRuSPawap\_w0lCtMA}): identificador interno único que Kafka asigna al topic; distingue al topic incluso si en el futuro se borrara y se recreara con el mismo nombre.
			\item \texttt{PartitionCount: 3}: el topic quedó dividido en 3 particiones, tal como se definió en \texttt{topics.yaml}. Cada partición es un log independiente y ordenado; dividir el topic en particiones es lo que permite paralelizar la escritura y la lectura entre varios brokers (y, más adelante, entre varios consumidores de un mismo grupo).
			\item \texttt{ReplicationFactor: 3}: cada partición tiene 3 copias (réplicas) distribuidas en los 3 brokers del clúster. Con exactamente 3 brokers y factor de replicación 3, cada broker termina teniendo una copia de \textbf{todas} las particiones (el máximo posible en este clúster), lo que da tolerancia a fallas: el clúster sigue funcionando aunque se caiga un broker.
			\item \texttt{Configs}: muestra la configuración efectiva del topic, incluyendo los overrides declarados en \texttt{topics.yaml} (\texttt{retention.ms=3600000}, es decir, los mensajes se retienen 1 hora; \texttt{cleanup.policy=delete}, se borran al vencer la retención) y valores por defecto del clúster (\texttt{min.insync.replicas=2}).
			\item Por cada partición (0, 1 y 2) se muestra su \texttt{Leader}: el broker que en este momento atiende las lecturas y escrituras de esa partición en particular. Nótese que el liderazgo quedó repartido entre los tres brokers (2, 3 y 1 respectivamente) en vez de concentrarse en uno solo, lo cual es deseable para balancear la carga.
			\item \texttt{Replicas} lista los brokers que tienen una copia de esa partición, e \texttt{Isr} (\textit{In-Sync Replicas}) el subconjunto de esas réplicas que está completamente al día con el líder. Que \texttt{Isr} sea idéntico a \texttt{Replicas} en las tres particiones indica que la replicación está sana, sin retraso.
			\item \texttt{Elr} y \texttt{LastKnownElr} (\textit{Eligible Leader Replicas}) son campos más recientes de Kafka relacionados con la elección segura de líder cuando el ISR se reduce por debajo de lo mínimo; aparecen vacíos porque no ha ocurrido ningún escenario de falla, lo cual es el estado esperado y saludable.
		\end{itemize}

		La Figura \ref{img:broker-partition-leader} resume visualmente esta misma información: qué broker es líder de cada partición y cuáles la replican como followers.

		\setindexcaption{Distribución de réplicas y liderazgo por partición del topic ecommerce-events.}
		\insertimage[\label{img:broker-partition-leader}]{kafka/broker-partition-leader.png}{width=14cm}{Cada broker guarda una réplica de las 3 particiones (factor de replicación 3), pero el liderazgo (en ámbar) queda repartido: Broker 1 lidera la Partición 2, Broker 2 lidera la Partición 0 y Broker 3 lidera la Partición 1.}
